% SYNTHETIC EVALUATION DRAFT — written for the retrieval eval in evals/.
% Not a real paper; contains no experimental results.
\documentclass{article}
\title{Few-Step Sampling for Latent Diffusion Models via Consistency Distillation}
\begin{document}
\maketitle

\begin{abstract}
Denoising diffusion probabilistic models generate high-fidelity images by
reversing a gradual noising process, but sampling requires tens to hundreds of
sequential network evaluations. This draft studies consistency distillation
for text-to-image latent diffusion models, aiming to produce images in one to
four denoising steps while preserving sample quality and prompt alignment
under classifier-free guidance.
\end{abstract}

\section{Introduction}
Diffusion models define a forward process that adds Gaussian noise to data over
many timesteps and train a neural network, typically a U-Net or a diffusion
transformer, to predict the noise (or the clean sample) at each step. Samples
are generated by starting from pure noise and iteratively denoising. Viewed in
continuous time, sampling corresponds to solving a reverse-time stochastic
differential equation or the associated probability flow ordinary
differential equation (ODE).

Latent diffusion models run this process in the latent space of a pretrained
variational autoencoder, which reduces cost and underlies most text-to-image
systems. Even so, high quality generation typically needs 25 to 50 solver
steps, and classifier-free guidance doubles the number of network evaluations
per step. This makes interactive and on-device generation difficult.

Our contributions are: (1) a guidance-aware consistency distillation objective
that distils a guided teacher into a student that maps any point on a
probability flow ODE trajectory directly to its origin; (2) a skipping-step
schedule that reduces discretization error of the teacher solver during
distillation; and (3) an evaluation plan covering image fidelity, diversity and
text alignment as a function of the number of sampling steps.

\section{Related Work}
DDPM introduced the denoising objective with a fixed variance schedule, and
score-based generative modelling unified diffusion with stochastic
differential equations. DDIM provided deterministic non-Markovian sampling with
fewer steps. Dedicated ODE solvers such as DPM-Solver and higher-order
exponential integrators further reduce the step count. Progressive
distillation repeatedly halves the number of steps of a student. Consistency
models learn a function that is self-consistent along ODE trajectories, and
latent consistency models apply the idea to Stable Diffusion. Adversarial
diffusion distillation and rectified flow are alternative routes to few-step
generation. Classifier-free guidance trades diversity for fidelity by
extrapolating between conditional and unconditional predictions.

\section{Method}
\subsection{Guidance-Aware Consistency Distillation}
Let $f_\theta(x_t, t, c)$ be the student consistency function conditioned on a
text embedding $c$. For a pair of adjacent timesteps $t_{n+k} > t_n$, we run one
step of the teacher ODE solver with classifier-free guidance scale $w$ from
$x_{t_{n+k}}$ to obtain $\hat{x}_{t_n}$, and minimize the distance between
$f_\theta(x_{t_{n+k}}, t_{n+k}, c)$ and an exponential moving average target
$f_{\theta^-}(\hat{x}_{t_n}, t_n, c)$. The guidance scale is sampled per batch
and provided to the student as an extra input, so a single student supports a
range of guidance strengths.

\subsection{Skipping-Step Schedule}
Using a skip interval $k > 1$ reduces the number of distillation targets per
trajectory and the accumulated teacher error. We plan to anneal $k$ during
training.

\section{Experimental Plan}
We plan to distil a pretrained text-to-image latent diffusion model and
evaluate samples at 1, 2, 4 and 8 steps. Metrics will be Fréchet Inception
Distance (FID), CLIP score for prompt alignment, and human preference ratings.
Baselines include DDIM, DPM-Solver, progressive distillation and latent
consistency models at matched step counts.

\section{Limitations}
Distilled students inherit the biases of the teacher model. One-step samples
tend to lose fine texture, and quality may degrade for prompts far from the
distillation prompt distribution.

\end{document}
